\chapter{Đào sâu: Gọi đồng bộ --- OpenFeign, Resilience4j và lỗi dây chuyền}
\label{ch:dd-resilience}

Chương~6 vạch ranh giới giữa hai cách service nói chuyện: HTTP khi bên
gọi cần câu trả lời ngay, sự kiện khi nó cần một sự thật được ghi
nhận. Chương này ở lại phía HTTP và hỏi câu hỏi mà mọi hệ thống phân
tán rồi cũng hỏi lúc 3 giờ sáng: khi service tôi đang gọi chậm, chuyện
gì xảy ra với \emph{tôi}? Dự án có đúng thư viện trên classpath ---
OpenFeign cho lời gọi, Resilience4j cho sự bảo vệ --- và, như bạn sẽ
thấy, đã nối chúng theo cách không bảo vệ được gì. Điều đó khiến nó
thành ca dạy học tốt hơn một ca hoàn chỉnh.

Ba câu hỏi sắp xếp chương này:

\begin{enumerate}
  \item \textbf{Một lời gọi đồng bộ tốn gì} ngoài độ trễ của nó --- và
  vì sao một phụ thuộc chậm tệ hơn một phụ thuộc chết?
  \item \textbf{Circuit breaker quyết định thế nào}, với các con số
  thật của dự án, và timeout, retry và breaker tương tác ra sao?
  \item \textbf{Fallback nên trả về gì}, và khi nào lời gọi không nên
  tồn tại?
\end{enumerate}

\section{Vì sao Feign, và nó giải quyết gì}

course-service thỉnh thoảng cần biết một user là ai --- một cái tên để
in lên bài nộp, một vai trò để kiểm tra. auth-service sở hữu dữ liệu
đó. Các lựa chọn là: gọi auth-service qua HTTP, chép dữ liệu sang
course-service qua sự kiện, hoặc đặt cả hai vào một cơ sở dữ liệu. Dự
án giữ các cơ sở dữ liệu tách biệt (Chương~7) và dành sự kiện cho các
sự thật đã xảy ra; tra cứu theo id là một câu hỏi, nên nó là một lời
gọi.

OpenFeign biến lời gọi đó thành một interface:

\begin{lstlisting}[language=Java]
@FeignClient(name = "auth-service",
             fallback = AuthServiceFallback.class)          (*@\co{1}@*)
public interface AuthServiceClient {

    @GetMapping("/users/{id}")                              (*@\co{2}@*)
    UserDto getUserById(@PathVariable("id") Long id);        (*@\co{3}@*)
}
\end{lstlisting}

\begin{description}
  \item[\co{1}] \code{name} là tên ứng dụng trên Eureka. Feign hỏi
  discovery client các instance của \code{auth-service} và cân bằng
  tải giữa chúng ở phía client --- không gateway, không host cứng.
  \code{fallback} nêu tên lớp sẽ được gọi thay khi lời gọi thất bại.
  Giữ ý đó lại.
  \item[\co{2}] Cùng các annotation Spring MVC như một controller, đọc
  ngược lại: đây là thứ cần \emph{gửi}. Path là của chính service
  (\code{/users/\{id\}}), không phải của gateway
  (\code{/api/auth/users}), vì lời gọi service-tới-service bỏ qua
  gateway.
  \item[\co{3}] Kiểu trả về là một record từ \code{common}. Record đó
  là \emph{hợp đồng}: auth-service phải sinh JSON giải tuần tự được
  thành nó.
\end{description}

Thứ Feign giải quyết là mã lặp và sự ràng buộc vào địa chỉ. Thứ nó
không giải quyết là tính chất duy nhất quan trọng dưới tải: mỗi lời gọi
giữ một thread Tomcat trong course-service lâu bằng thời gian
auth-service cần để trả lời.

\begin{decision}[một lời gọi, không phải một bản sao --- tạm thời]
Với một tra cứu theo id hiếm khi xảy ra và câu trả lời hiếm khi đổi,
một lời gọi là thứ đúng đơn giản nhất, và Feign làm nó trong hai dòng.
Quyết định sẽ lật khi bất kỳ điều nào trong ba điều thành sự thật: tra
cứu nằm trên đường nóng (mọi danh sách khóa học cần tên giáo viên), dữ
liệu cần có trong lúc auth-service chết (chấm bài không được dừng vì
đăng nhập hỏng), hoặc hình dạng của \code{UserDto} bắt đầu đổi theo các
bản phát hành của auth-service. Khi đó course-service nên giữ bản sao
riêng của những trường nó cần, nuôi bằng các sự kiện auth-service đã
phát (\code{ts.user.registered}, \code{ts.user.activated}) --- mẫu mà
Chương~7 đã áp dụng cho học sinh.
\end{decision}

\section{Hợp đồng trong \code{common}, và rủi ro phiên bản của nó}

\code{UserDto} là một record với năm trường: \code{id}, \code{email},
\code{firstName}, \code{lastName}, \code{role}. \code{UserController}
của auth-service không trả về nó --- nó trả về \code{UserResponse} của
riêng mình. Hai lớp hôm nay khớp nhau nhờ trùng hợp tên trường, và mặc
định bỏ qua thuộc tính lạ của Jackson che đi khoảng hở: auth-service
có thể thêm trường thoải mái, và course-service sẽ không nhận ra. Nó
\emph{sẽ} nhận ra một lần đổi tên. Đổi \code{firstName} thành
\code{givenName} trong \code{UserResponse} biên dịch được ở cả hai
phía, deploy được ở cả hai phía, và để course-service đọc tên là
\code{null} --- không ngoại lệ, không dòng log. Phiên bản phỏng vấn:
``bạn đánh phiên bản một API nội bộ thế nào?'' và câu trả lời là: chỉ
thay đổi có tính bổ sung, không bao giờ đổi tên hay xóa; DTO của bên
tiêu thụ bỏ qua thuộc tính lạ; một consumer-driven contract test
(Spring Cloud Contract, Pact) trong build của auth-service thất bại khi
\code{UserResponse} không còn thỏa \code{UserDto}. Phiên bản trong URL
(\code{/v2/users}) dành cho thay đổi phá vỡ không tránh được, và nó
nghĩa là chạy cả hai chừng nào còn bên tiêu thụ nào chưa theo kịp.

\section{Circuit breaker, với các con số của dự án}
\label{sec:cb-math}

Breaker của Resilience4j là một máy trạng thái trên một \emph{cửa sổ
trượt} các kết quả gọi gần đây. Cấu hình của dự án trong
\code{course-service.yml}:

\begin{lstlisting}[language=yaml]
resilience4j:
  circuitbreaker:
    instances:
      authService:
        sliding-window-size: 10                    (*@\co{1}@*)
        failure-rate-threshold: 50                 (*@\co{2}@*)
        wait-duration-in-open-state: 10s           (*@\co{3}@*)
        permitted-number-of-calls-in-half-open-state: 3  (*@\co{4}@*)
\end{lstlisting}

\begin{figure}[htb]
\centering
\begin{tikzpicture}[node distance=13mm and 14mm]
  \node[boxfill] (closed) {CLOSED\\lời gọi đi qua, kết quả được ghi};
  \node[boxdark, right=of closed] (open) {OPEN\\từ chối tức thì};
  \node[box, below=of $(closed)!0.5!(open)$] (half) {HALF\_OPEN\\3 lời gọi thử};
  \draw[flow] (closed) -- node[lbl,above]{$\ge$5 trong 10 lỗi} (open);
  \draw[flow] (open.east) -- ++(6mm,0) |- node[lbl,above,pos=0.75]{sau 10\,s} (half.east);
  \draw[flow] (half.north west) -- node[lbl,above,sloped]{thử ổn} (closed.south);
  \draw[flowdash] (half.north east) -- node[lbl,above,sloped]{thử lỗi} (open.south);
\end{tikzpicture}
\caption{Ba trạng thái của breaker với các ngưỡng của dự án này. Lời
gọi bị từ chối ở OPEN đi tới fallback mà không chạm mạng.}
\label{fig:cb-states}
\end{figure}

\begin{description}
  \item[\co{1}] Cửa sổ giữ 10 kết quả gần nhất (theo số lượng; phương
  án kia là theo thời gian, ``60 giây gần nhất''). Breaker không đánh
  giá cho tới khi cửa sổ có đủ \code{minimum-number-of-calls} kết quả
  --- không đặt ở đây, nên mặc định bằng cỡ cửa sổ. Dịch ra: 9 lời gọi
  đầu sau khởi động có thể \emph{toàn bộ} thất bại mà breaker vẫn
  CLOSED.
  \item[\co{2}] Ở 10 kết quả, nếu 5 trở lên thất bại, breaker mở. Một
  ``thất bại'' là một ngoại lệ breaker không bỏ qua. Để ý thứ
  \emph{không} được cấu hình: \code{slow-call-duration-threshold}. Mặc
  định của nó là 60\,s, nên một lời gọi mất 30\,s là một \emph{thành
  công} với breaker này. Phụ thuộc chậm --- trường hợp phổ biến ---
  không bao giờ làm nó nhảy.
  \item[\co{3}] OPEN trong 10\,s: mọi lời gọi bị từ chối bằng
  \code{CallNotPermittedException} trong micro giây và chuyển tới
  fallback. auth-service được nghỉ 10 giây.
  \item[\co{4}] Rồi HALF\_OPEN: 3 lời gọi được cho qua làm thăm dò. Nếu
  tỷ lệ thất bại của chúng $\ge$ 50\,\% (2 trong 3), quay lại OPEN thêm
  10\,s; nếu không, CLOSED với cửa sổ mới.
\end{description}

Giờ tới phần khiến chương này thẳng thắn. Instance \code{authService}
đó không gắn với \emph{gì cả}. Không phương thức nào mang
\code{@CircuitBreaker(name = "authService")}, và tích hợp riêng của
Feign --- thứ sẽ bọc \code{getUserById} trong một breaker tên
\code{AuthServiceClient\#getUserById(Long)} --- đang tắt, vì công tắc
\code{openfeign.circuitbreaker.enabled} mặc định là \code{false}, và
starter cung cấp \code{CircuitBreakerFactory} nó cần (starter
\code{circuitbreaker-resilience4j}) không có trên classpath. Với công
tắc đó tắt, thuộc tính \code{fallback} ở \co{1} bị lặng lẽ bỏ qua. Và
\code{AuthServiceClient} hôm nay không được gọi từ lớp service nào. Dự
án có một breaker cấu hình cho một lời gọi không xảy ra, bảo vệ không
gì cả, với một fallback sẽ không chạy. Đó là mẫu vật hoàn hảo của cách
khả năng chịu lỗi thường bị cấu hình sai trong production: mọi mảnh có
mặt, không mảnh nào được nối.

\section{Timeout xếp chồng}

Khi lời gọi được nối đúng, ba bộ đếm giờ bao quanh nó, và chúng phải
được sắp thứ tự nếu không các bộ bên ngoài chỉ là trang trí.

\begin{itemize}
  \item \textbf{Timeout socket của Feign.} \code{connectTimeout} và
  \code{readTimeout} trên HTTP client. Không đặt trong dự án, nên mặc
  định áp dụng: 10\,s để kết nối, \textbf{60\,s} để đọc. Một
  auth-service treo ghim một thread course-service suốt một phút.
  \item \textbf{TimeLimiter của Resilience4j.} Một hạn chót cho toàn
  bộ lời gọi được bọc, độc lập với socket. Nó chỉ áp dụng cho kiểu trả
  về bất đồng bộ (\code{CompletableFuture}), đó là lý do lời gọi Feign
  đồng bộ phải dựa vào timeout socket.
  \item \textbf{Hạn chót của bên gọi.} \code{response-timeout} của
  gateway (Chương~31), và phía sau nó là sự kiên nhẫn của trình duyệt.
  Nếu gateway bỏ cuộc ở 5\,s trong khi Feign đợi 60\,s, người dùng thấy
  lỗi ở 5\,s và course-service giữ một thread bận thêm 55 giây để phục
  vụ không ai.
\end{itemize}

Quy tắc: timeout của mỗi chặng \emph{ngắn hơn} chặng trước nó. Trình
duyệt 10\,s $>$ gateway 5\,s $>$ Feign read của course-service 2\,s.
Một timeout nổ bên trong một request đã bị bỏ rơi là công sức lãng
phí; một timeout nổ muộn hơn của bên gọi là một chỗ rò.

\section{Retry: chỉ cho lời gọi lặp lại được}

\code{@Retry} của Resilience4j phát lại một lời gọi thất bại. Trên
\code{GET /users/\{id\}} thì vô hại: đọc hai lần cũng như đọc một lần.
Trên \code{POST /enrollments} thì là một lần ghi danh trùng, vì timeout
không cho bạn biết request đầu đã được xử lý hay chưa. Vì thế retry chỉ
dành cho thao tác \emph{lũy đẳng} (idempotent), hoặc thao tác được làm
cho lũy đẳng bằng một khóa client cung cấp mà server khử trùng lặp
theo đó. Thêm hai quy tắc từ production: retry với \emph{lùi theo hàm
mũ có jitter}, nếu không mọi client retry cùng một khoảnh khắc và đợt
sóng thứ hai lớn hơn đợt đầu (``bão retry''); và retry \emph{bên trong}
breaker, để một breaker đang mở chặn retry ngay từ đầu. Thứ tự aspect
mặc định của Resilience4j --- Retry ngoài cùng, rồi CircuitBreaker,
rồi TimeLimiter, rồi Bulkhead --- làm đúng điều này: mỗi lần retry đi
qua breaker, và lời gọi bị từ chối thì không được retry.

\section{Bulkhead: giới hạn vùng nổ}

Bulkhead giới hạn số lời gọi đồng thời được phép đang bay tới một phụ
thuộc. Bulkhead semaphore: một bộ đếm; lời gọi đồng thời thứ 26 bị từ
chối ngay. Bulkhead thread pool: lời gọi chạy trên một pool riêng để
thread của bên gọi không bao giờ là kẻ đứng chờ --- cần thiết với một
ứng dụng servlet có pool Tomcat hữu hạn. course-service chạy Tomcat với
200 thread. Không có bulkhead, 200 lời gọi auth-service chậm tiêu sạch
từng thread, và course-service không trả lời nổi một probe sức khỏe,
nói gì tới danh sách khóa học. Với bulkhead semaphore 25, tối đa 25
thread từng chờ auth-service và 175 còn lại cho mọi thứ khác. Đó là sự
cô lập mà tiêu đề chương hứa hẹn.

\section{Fallback trả về gì}

\begin{lstlisting}[language=Java]
@Component
public class AuthServiceFallback implements AuthServiceClient {
    @Override
    public UserDto getUserById(Long id) {
        return null;                                       (*@\co{1}@*)
    }
}
\end{lstlisting}

\begin{description}
  \item[\co{1}] Một \code{null} mà bên gọi phải nhớ kiểm tra. Mọi nơi
  tiêu thụ \code{getUserById} giờ có một nhánh ẩn, và một lần quên kiểm
  tra là một \code{NullPointerException} chỉ nổ khi auth-service chết
  --- đúng lúc bạn ít muốn một lỗi mới nhất.
\end{description}

Một fallback có bốn lựa chọn trung thực, theo thứ tự ưu tiên: một giá
trị \emph{đã cache} từ lời gọi thành công gần nhất (cũ nhưng thật);
một giá trị \emph{suy giảm} mà UI hiển thị được (``Giáo viên không
rõ'') và được định kiểu rõ là như vậy; một kết quả \emph{rỗng} nơi
miền cho phép; hoặc một \emph{lỗi có kiểu} mà bên gọi ánh xạ được thành
503 với \code{Retry-After}. \code{null} không phải cái nào trong số đó.
Nó là sự vắng mặt của một quyết định, đẩy sang mọi bên gọi. Fallback
cũng cần \emph{nguyên nhân}: một phương thức fallback hai tham số
(\code{Throwable last}) phân biệt được một lần breaker từ chối với một
404, và 404 --- user không tồn tại --- \emph{không} được biến thành
``thử lại sau''.

\section{Ưu và nhược, thẳng thắn}

\begin{center}
\small
\begin{tabular}{@{}p{0.46\linewidth}p{0.46\linewidth}@{}}
\toprule
Dự án có gì & Thiếu gì \\
\midrule
Client khai báo, phân giải qua Eureka & Client không bao giờ được gọi \\
Một instance breaker với ngưỡng hợp lý & Không gì gắn vào nó; công tắc breaker của Feign đang tắt \\
Một lớp fallback đã đăng ký & Bị bỏ qua; và nó trả về \code{null} \\
Resilience4j trên classpath & Không timeout, không retry, không bulkhead, không ngưỡng chậm \\
Actuator đã mở & Endpoint breaker không nằm trong danh sách mở \\
\bottomrule
\end{tabular}
\end{center}

Và thêm một khoảng hở không liên quan tới khả năng chịu lỗi:
\code{SecurityConfig} của auth-service chỉ cho phép các đường dẫn công
khai; \code{/users/\{id\}} cần một ngữ cảnh đã xác thực, thứ
\code{GatewayHeaderAuthFilter} của nó dựng từ \code{X-User-Id}. Feign
không gửi header nào như vậy. Lời gọi, một khi được nối, sẽ trả 401 ở
mọi lần thử --- và 401 là một ngoại lệ, nên mười lần sẽ mở breaker
chống lại một service hoàn toàn khỏe mạnh.

\section{Cái gì gãy}

\begin{pitfall}[fallback không bao giờ chạy]
Triệu chứng: auth-service bị dừng, và thay vì \code{null} của fallback,
course-service ném \code{FeignException} với stack trace xuyên qua load
balancer, và request trả về 500. Nguyên nhân:
\code{@FeignClient(fallback = ...)} chỉ được tôn trọng khi
\code{spring.cloud.openfeign.circuitbreaker.enabled=true} \emph{và} có
một bean \code{CircuitBreakerFactory}, thứ cần dependency
\code{spring-cloud-starter-circuitbreaker-resilience4j}. Thiếu một
trong hai, Spring không log gì và thuộc tính trở nên trơ. Sửa ở
Mục~\ref{sec:res-enhance}.
\end{pitfall}

\begin{pitfall}[service khỏe, breaker mở]
Triệu chứng: sau một lần deploy auth-service, breaker của course-service
mở và cứ mở, quay vòng OPEN $\to$ HALF\_OPEN $\to$ OPEN, dù
auth-service trả lời \code{curl} bình thường. Nguyên nhân: mọi lời gọi
nhận 401 (thiếu header danh tính), hoặc 404 cho một user đã xóa, và cả
hai đều tính là thất bại. Sửa: lan truyền danh tính trên lời gọi
service-tới-service (một \code{RequestInterceptor}, bên dưới), và bảo
breaker ngoại lệ nào \emph{không} phải lỗi phụ thuộc bằng
\code{ignore-exceptions} --- lỗi phía client nghĩa là request sai, không
phải service chết.
\end{pitfall}

\begin{pitfall}[thread 60 giây]
Triệu chứng: auth-service treo (deadlock, GC kẹt); CPU của
course-service rảnh, vậy mà nó ngừng trả lời \emph{mọi} request, và
readiness probe của nó thất bại nên Kubernetes ngừng route tới nó.
Nguyên nhân: read timeout mặc định 60\,s của Feign nhân với 200 thread
của Tomcat. Mỗi tra cứu treo giữ một thread một phút; với tốc độ
request vừa phải, mọi thread đều đứng chờ trong vài giây. Breaker không
giúp được --- với ngưỡng chậm mặc định, lời gọi treo không phải thất
bại cho tới khi hết timeout. Sửa: read timeout 2\,s, một ngưỡng chậm,
và bulkhead 25.
\end{pitfall}

\section{Chuyện gì gãy lúc 3 giờ sáng}

Postgres của auth-service bắt đầu một đợt vacuum dài;
\code{/users/\{id\}} đi từ 20\,ms lên 2\,s. Giả sử lời gọi Feign đã
được nối vào endpoint danh sách khóa học, như thiết kế của dự án mời
gọi, và không gì khác thay đổi.

\begin{enumerate}
  \item Mỗi danh sách khóa học giờ mất 2\,s thay vì 20\,ms. Thread
  Tomcat trong course-service bận lâu hơn 100$\times$ cho mỗi request.
  \item Ở 50 danh sách mỗi giây, 100 thread đang chờ auth-service tại
  bất kỳ thời điểm nào. Ở 100 mỗi giây, cả 200 đều chờ, và request xếp
  hàng trong backlog accept của Tomcat.
  \item Readiness probe của course-service (\code{/actuator/health}
  thậm chí không được mở ở đây --- Chương~17 dùng probe TCP) vẫn qua,
  vì socket vẫn nhận kết nối. Gateway tiếp tục gửi lưu lượng.
  \item Connection pool của gateway tới course-service đầy; route
  \code{/api/courses} của nó suy giảm cho tất cả mọi người, kể cả các
  request chưa từng cần auth-service. Chuỗi lỗi của Chương~31 bắt đầu
  cao hơn một chặng.
  \item auth-service, trong khi đó, nhận 100 request mỗi giây từ một
  bên gọi sẽ vứt câu trả lời sau timeout của gateway --- tải khiến sự
  tranh chấp của đợt vacuum tệ hơn.
  \item Breaker không bao giờ mở: 2\,s dưới mặc định chậm 60\,s, và
  không lời gọi nào thất bại. Cơ chế duy nhất có thể cắt tải đang nhìn
  vào sai tín hiệu.
\end{enumerate}

Với cấu hình bên dưới: lời gọi bị cắt ở 2\,s (Feign read timeout) và
tính là chậm trên 1\,s; sau mười kết quả với $\ge$80\,\% chậm, breaker
mở; danh sách trả về với tên giáo viên đã cache hoặc suy giảm dưới một
mili giây; tối đa 25 thread từng đứng chờ; tải của auth-service giảm
còn ba lời gọi thăm dò mỗi mười giây; và thay đổi trạng thái của
breaker nằm trên một dashboard.

\section{Nâng cấp giải pháp}
\label{sec:res-enhance}

Dependency trước: thêm \code{spring-cloud-starter-circuitbreaker-resilience4j}
vào POM của course-service (nó mang tới \code{CircuitBreakerFactory} mà
Feign cần; \code{resilience4j-spring-boot3} giữ lại cho annotation và
metric). Rồi \code{course-service.yml}:

\begin{lstlisting}[language=yaml]
spring:
  cloud:
    openfeign:
      circuitbreaker:
        enabled: true                            (*@\co{1}@*)
      client:
        config:
          auth-service:                          (*@\co{2}@*)
            connect-timeout: 1000
            read-timeout: 2000
            logger-level: basic

resilience4j:
  circuitbreaker:
    instances:
      authService:
        sliding-window-type: COUNT_BASED
        sliding-window-size: 10
        minimum-number-of-calls: 10
        failure-rate-threshold: 50
        slow-call-duration-threshold: 1s         (*@\co{3}@*)
        slow-call-rate-threshold: 80
        wait-duration-in-open-state: 10s
        permitted-number-of-calls-in-half-open-state: 3
        automatic-transition-from-open-to-half-open-enabled: true
        register-health-indicator: true          (*@\co{4}@*)
        ignore-exceptions:                       (*@\co{5}@*)
          - feign.FeignException.NotFound
          - feign.FeignException.BadRequest
  retry:
    instances:
      authService:
        max-attempts: 3                          (*@\co{6}@*)
        wait-duration: 200ms
        enable-exponential-backoff: true
        exponential-backoff-multiplier: 2
        enable-randomized-wait: true
        randomized-wait-factor: 0.5
        retry-exceptions:
          - java.net.SocketTimeoutException
          - feign.RetryableException
        ignore-exceptions:
          - feign.FeignException.FeignClientException
  bulkhead:
    instances:
      authService:
        max-concurrent-calls: 25                 (*@\co{7}@*)
        max-wait-duration: 0
  timelimiter:
    instances:
      authService:
        timeout-duration: 3s                     (*@\co{8}@*)
        cancel-running-future: true

management:
  endpoints:
    web:
      exposure:
        include: health,info,metrics,prometheus,circuitbreakers,circuitbreakerevents
  health:
    circuitbreakers:
      enabled: true
\end{lstlisting}

\begin{description}
  \item[\co{1}] Công tắc khiến \code{fallback = ...} thành thật. Feign
  giờ bọc mỗi phương thức trong một breaker tên
  \code{AuthServiceClient\#getUserById(Long)}; các annotation tường
  minh bên dưới dùng instance \code{authService} cho chính sách phía
  bên gọi.
  \item[\co{2}] Timeout theo từng client, khóa theo \code{name} của
  Feign. Hai giây là rộng rãi cho một tra cứu theo id bình thường mất
  20\,ms, và ngắn hơn hẳn 5\,s của gateway.
  \item[\co{3}] Chậm giờ được tính. Trên 1\,s là ``chậm hơn thiết kế'';
  khi 80\,\% cửa sổ là chậm, breaker mở dù mọi lời gọi rồi cũng thành
  công.
  \item[\co{4}] Trạng thái breaker xuất hiện trong
  \code{/actuator/health} và dưới dạng một gauge Micrometer
  (\code{circuitbreaker.state}) --- thứ bạn đặt cảnh báo.
  \item[\co{5}] Một user không tồn tại là một câu trả lời, không phải
  sự cố. Phản hồi 4xx không tính là thất bại cũng không kích hoạt
  fallback.
  \item[\co{6}] Ba lần thử với 200\,ms, rồi 400\,ms, mỗi lần jitter
  $\pm$50\,\%. Chỉ lỗi vận chuyển được retry; bất kỳ HTTP 4xx nào là
  chung cuộc. Tổng trường hợp xấu nhất:
  $3 \times 2\,\text{s} + 0{,}6\,\text{s}$, vẫn có giới hạn và vẫn dưới
  một phút --- nhưng dài hơn 5\,s của gateway, đó là lý lẽ cho tối đa 2
  lần thử trên đường request và 3 chỉ trong công việc nền.
  \item[\co{7}] Hai mươi lăm thread được phép chờ auth-service; lời gọi
  thứ 26 thất bại ngay với \code{BulkheadFullException} và đi
  fallback. \code{max-wait-duration: 0} nghĩa là không xếp hàng --- xếp
  hàng chính là thứ ta đang tránh.
  \item[\co{8}] Chỉ dùng bởi biến thể \code{CompletableFuture} bên
  dưới; vô hại nếu không.
\end{description}

Lớp bọc phía service, với lan truyền danh tính và một fallback nói rõ
nó là gì:

\begin{lstlisting}[language=Java]
/** Forwards the caller's identity on service-to-service calls so
 *  auth-service's GatewayHeaderAuthFilter accepts them. */
@Configuration
public class FeignIdentityConfig {
    @Bean
    public RequestInterceptor identityPropagation() {
        return template -> {
            var auth = SecurityContextHolder.getContext()
                    .getAuthentication();
            if (auth != null && auth.getPrincipal() instanceof Long id) {
                template.header(JwtConstants.HEADER_USER_ID, id.toString());
                auth.getAuthorities().stream().findFirst().ifPresent(a ->
                    template.header(JwtConstants.HEADER_USER_ROLE,
                        a.getAuthority().replace("ROLE_", "")));
            }
        };
    }
}

/** The one place course-service asks who a user is. */
@Service
@RequiredArgsConstructor
public class UserLookupService {

    private final AuthServiceClient authServiceClient;
    private final Map<Long, UserDto> lastKnown =
            new ConcurrentHashMap<>();                     (*@\co{1}@*)

    @Retry(name = "authService")
    @CircuitBreaker(name = "authService", fallbackMethod = "degraded")
    @Bulkhead(name = "authService")
    public UserDto findUser(Long id) {
        UserDto user = authServiceClient.getUserById(id);
        lastKnown.put(id, user);
        return user;
    }

    /** Same signature plus the cause; called for breaker rejections,
     *  bulkhead rejections, timeouts and 5xx -- never for 4xx. */
    private UserDto degraded(Long id, Throwable cause) {   (*@\co{2}@*)
        UserDto cached = lastKnown.get(id);
        if (cached != null) {
            return cached;                                 // stale, real
        }
        return new UserDto(id, null, "Unknown", "user", null);
    }
}
\end{lstlisting}

\begin{description}
  \item[\co{1}] Một cache giá-trị-tốt-gần-nhất. Với một replica, một
  map là đủ; với nhiều replica, cùng ý tưởng đó thuộc về Redis với TTL,
  hoặc --- cách sửa thật --- một bảng cục bộ nuôi bằng sự kiện user.
  \item[\co{2}] Fallback có kiểu và nhìn thấy được: bên gọi nhận một
  \code{UserDto} có tên nói ``Unknown'', không bao giờ \code{null}.
  \code{Throwable} cho bạn log \emph{vì sao} (breaker mở hay timeout)
  ở mức WARN một lần mỗi lần đổi trạng thái thay vì mỗi lời gọi.
\end{description}

Các metric cần theo dõi, tất cả đã xuất sang Prometheus bởi thiết lập
Actuator sẵn có: \code{resilience4j\_circuitbreaker\_state} (một gauge
theo từng trạng thái từng instance; cảnh báo khi \code{state="open"}
bằng 1 quá một phút), \code{resilience4j\_circuitbreaker\_calls\_seconds}
theo \code{kind} (successful, failed, ignored) cho tốc độ,
\code{resilience4j\_bulkhead\_available\_concurrent\_calls} tiến về
không như chỉ báo sớm, và \code{http\_client\_requests\_seconds} của
chính Feign cho độ trễ thô đã khởi đầu tất cả.

\begin{decision}[khi nào lời gọi không nên tồn tại]
Giữ lời gọi đồng bộ chừng nào nó nằm ngoài đường nóng, hiếm, và chịu
được câu trả lời cũ. Thay nó bằng bản sao cục bộ nuôi bằng sự kiện khi
một trang danh sách cần nó (độ trễ và ràng buộc), khi việc chấm bài
phải chạy trong lúc auth-service chết (tính sẵn sàng), hoặc khi hai đội
phát hành theo nhịp khác nhau (trôi hợp đồng). Thay nó bằng một cache
TTL ngắn đặt trước lời gọi khi dữ liệu nóng nhưng không được cũ quá một
phút và nhất quán cuối cùng qua sự kiện không đáng một consumer mới.
Bản năng senior là đếm số phụ thuộc đồng bộ trên mỗi request và coi mỗi
cái là một phép nhân xác suất lỗi: ba phụ thuộc 99,9\,\% nối tiếp là
một endpoint 99,7\,\%.
\end{decision}

\begin{handson}[xem breaker mở]
Với phần nâng cấp đã áp dụng và stack đang chạy, dừng auth-service và
gọi một endpoint dùng \code{UserLookupService} mười hai lần:
\begin{lstlisting}[language=bash]
$ docker compose stop auth-service
$ for i in $(seq 1 12); do
    curl -s -o /dev/null -w "%{http_code} %{time_total}s\n" \
      -H "Authorization: Bearer $TOKEN" \
      http://localhost:8080/api/courses/courses/1
  done
200 1.04s   ... (10 slow lines: retries then fallback)
200 0.01s   <- breaker open: fallback without a network call
200 0.01s
$ curl -s http://localhost:8082/actuator/circuitbreakers | jq
{ "circuitBreakers": { "authService": {
    "failureRate": "100.0%", "slowCallRate": "0.0%",
    "bufferedCalls": 10, "failedCalls": 10, "state": "OPEN" } } }
\end{lstlisting}
Mọi phản hồi đều là 200 với ``Unknown user'' làm giáo viên --- trang
suy giảm thay vì thất bại. Khởi động lại auth-service, đợi 10\,s, và
\code{/actuator/circuitbreakerevents} cho thấy ba lời gọi thăm dò
HALF\_OPEN thành công và trạng thái trở về CLOSED.
\end{handson}

\section{Ngoài dự án}

Cùng bộ máy dưới những cái tên khác: Hystrix là bản gốc của Netflix và
đã nghỉ hưu; Resilience4j là kẻ kế nhiệm trong thế giới Java; Envoy và
Istio hiện thực timeout, retry có ngân sách, và loại trừ outlier trong
sidecar để code ứng dụng không chứa gì trong số đó, hướng đi mà các đội
tàu lớn đã chọn. gRPC mang hạn chót (deadline) trong giao thức --- hạn
chót của bên gọi lan truyền tới mọi chặng, giải quyết bài toán
``timeout xếp chồng'' ngay từ cấu trúc. SDK cloud (AWS, GCP) đi kèm lùi
theo hàm mũ có jitter sẵn vì cùng lý do chương này khuyến nghị. Từ vựng
là phổ quát; một người phỏng vấn nói ``bulkhead'' hay ``retry budget''
có ý đúng như chương này.

\section{Câu hỏi từ phỏng vấn thật}

\subsection*{Q: Độ trễ auth-service đi từ 20\,ms lên 2\,s. Đi qua với tôi chuyện gì xảy ra trong course-service, và bạn chặn nó thế nào.}

Mỗi request chạm auth-service giờ giữ một thread Tomcat lâu gấp trăm
lần. Ở tốc độ vừa phải, pool 200 bão hòa trong vài giây, hàng đợi
accept đầy, và course-service ngừng trả lời mọi thứ --- kể cả các
endpoint chưa từng gọi auth-service và readiness probe. Ở thượng nguồn,
pool theo host của gateway tới course-service đầy và route của nó suy
giảm cho mọi user. Không gì trong cấu hình mặc định chặn được điều này:
Feign đợi 60\,s, và một lời gọi 2\,s không phải chậm với một breaker có
ngưỡng chậm 60\,s. Để chặn, theo thứ tự đòn bẩy: read timeout 2\,s để
không thread nào đợi lâu hơn hạn chót của gateway; bulkhead khoảng 25
để 175 thread còn lại tiếp tục phục vụ; ngưỡng chậm 1\,s ở 80\,\% để
breaker mở theo độ trễ, không chỉ theo lỗi; và một fallback trả về tên
đã biết gần nhất thay vì \code{null}. Rồi, nếu tra cứu nằm trên đường
nóng, bỏ hẳn lời gọi bằng bản sao cục bộ nuôi bằng sự kiện user. Cách
đặt vấn đề senior: mục tiêu là độ trễ của auth-service hiện ra dưới
dạng một cái tên cũ trên trang, không phải một sự cố của khóa học.

\subsection*{Q: Bạn định cỡ timeout qua ba chặng thế nào?}

Từ ngoài vào trong, mỗi cái ngắn hơn cái trước. Trình duyệt hay client
di động là ngân sách ngoài cùng --- chẳng hạn 10\,s trước khi người
dùng bỏ cuộc. Timeout phản hồi của gateway phải nằm trong đó, 5\,s, để
client nhận 503 sạch sẽ thay vì treo. Feign read timeout của
course-service phải nằm trong \emph{cái đó} và chừa chỗ cho retry:
2\,s với tối đa một lần retry giữ trường hợp xấu nhất quanh 4\,s. Bất
kỳ timeout nào dài hơn của bên gọi là một chỗ rò --- công sức bỏ ra cho
một request đã bị bỏ rơi --- và bất kỳ chính sách retry nào cũng phải
tính vào phép cộng, vì ba lần thử ở 2\,s là 6\,s cộng backoff. Trong
production bạn cũng đặt timeout \emph{kết nối} riêng và ngắn hơn nhiều
(1\,s), vì không kết nối được thì chẩn đoán trong một vòng đi-về. Các
hệ thống tốt nhất lan truyền một header hạn chót từ biên để mỗi chặng
tự tính ngân sách còn lại; gRPC làm điều này ngay trong giao thức.

\subsection*{Q: Bão retry --- bạn ngăn thế nào?}

Ba biện pháp, cần tất cả. Thứ nhất, jitter: lùi theo hàm mũ đơn thuần
để mọi client thất bại cùng lúc retry cùng lúc, nên đợt sóng thứ hai
đồng bộ và lớn hơn; ngẫu nhiên hóa thời gian chờ $\pm$50\,\% dàn nó
ra. Thứ hai, ngân sách retry: retry theo phần trăm lưu lượng (mặc định
của Envoy là 20\,\%), để khi một phụ thuộc thực sự chết, retry không
thể nhân ba tải của nó. Resilience4j xấp xỉ điều này bằng cách đặt
Retry ngoài breaker --- một breaker đang mở từ chối các lần thử trước
khi chúng thành retry. Thứ ba, tính lũy đẳng: chỉ retry \code{GET} và
các thao tác mang khóa lũy đẳng mà server khử trùng lặp; nếu không một
timeout trên \code{POST} thành một lần ghi danh trùng. Và chỉ retry lỗi
vận chuyển và 503 --- một 400 hay 404 là tất định và sẽ lỗi lại. Dấu
hiệu của bão trong metric là một phụ thuộc có tốc độ request tăng khi
tỷ lệ thành công của nó giảm.

\subsection*{Q: Breaker đang mở, nhưng user cần dữ liệu. Fallback của bạn trả về gì?}

Thứ trung thực nhất có sẵn, định kiểu để bên gọi không nhầm nó với giá
trị thật. Thứ tự ưu tiên: giá trị thành công gần nhất từ một cache sống
ngắn (cũ nhưng thật --- tên giáo viên của mười phút trước là ổn); một
giá trị thay thế suy giảm rõ ràng mà UI render được (``Giáo viên không
rõ'') không bao giờ giả dạng dữ liệu; một tập rỗng nơi miền cho phép;
hoặc một ngoại lệ có kiểu trở thành 503 với \code{Retry-After}. Không
bao giờ \code{null}, thứ biến một sự kiện chịu lỗi thành
\code{NullPointerException} ở đâu đó khác. Fallback cũng phải thấy
nguyên nhân: 404 nghĩa là user không tồn tại và nên hiện ra như thế,
không bị nuốt vào ``thử lại sau''. Và fallback phải rẻ và cục bộ ---
một fallback gọi service khác là một phụ thuộc thứ hai trên đường lỗi.

\subsection*{Q: Làm sao bạn biết breaker đang mở lúc 3 giờ sáng?}

Vì nó là một metric có cảnh báo, không phải một dòng log. Resilience4j
xuất \code{resilience4j\_circuitbreaker\_state} như một gauge theo từng
trạng thái; một rule như ``\code{state="open"} bằng 1 quá 60 giây''
gọi ai đó với tên phụ thuộc trong cảnh báo. Ghép với các chỉ báo sớm:
\code{slow\_call\_rate} tăng trước khi breaker nhảy, số chỗ trống
bulkhead tiến về không, và p99 của phụ thuộc từ
\code{http\_client\_requests\_seconds}. Breaker cũng nuôi
\code{/actuator/health} khi bật \code{register-health-indicator}, nên
readiness probe có thể phản ánh nó --- dù thường bạn \emph{không} muốn
một breaker mở kéo service ra khỏi vòng xoay, vì nó đang phục vụ
fallback chính là để ở lại. Và log \emph{lần chuyển trạng thái} ở WARN
một lần, không phải mỗi lời gọi bị từ chối, nếu không log thành sự cố
kế tiếp.

\subsection*{Q: Bạn đặt breaker ở gateway hay ở service?}

Cả hai, cho hai loại lỗi khác nhau. Breaker của gateway (Chương~31)
bảo vệ user khỏi cả một service chậm; nó mở theo route và phục vụ một
body fallback. Breaker của service bảo vệ thread pool của chính service
khỏi một trong các phụ thuộc của nó và là nơi fallback hiểu miền sống
--- chỉ course-service biết một tên giáo viên thiếu là sống được.
Gateway không thể biết điều đó. Sai lầm là chỉ có cái ở biên và cho
rằng các service phía sau vì thế an toàn; chuỗi lỗi mà chương này đi
qua diễn ra hoàn toàn phía sau gateway.

\begin{quiz}
\begin{enumerate}
  \item Dự án cấu hình một breaker tên \code{authService}. Liệt kê ba
  lý do riêng biệt khiến nó hiện không bảo vệ lời gọi nào, và thay đổi
  một dòng sửa từng lý do.
  \item Với \code{sliding-window-size: 10} và các mặc định còn lại, bao
  nhiêu lần thất bại liên tiếp có thể xảy ra sau khởi động trước khi
  breaker mở, và vì sao một lời gọi mất 30\,s không được tính?
  \item Sắp thứ tự và biện minh: timeout phản hồi của gateway, timeout
  của trình duyệt, Feign read timeout, số lần retry.
  \item \code{AuthServiceFallback.getUserById} trả về \code{null}. Nêu
  hai khiếm khuyết cụ thể kéo theo và chữ ký hai tham số sửa được cái
  tinh vi hơn.
  \item auth-service trả 401 cho mọi lời gọi Feign. Giải thích nguyên
  nhân trong cấu hình bảo mật của dự án, tác động lên breaker, và hai
  thay đổi giải quyết nó.
  \item Ba phụ thuộc đồng bộ mỗi cái 99,9\,\% sẵn sàng nối tiếp trên
  một endpoint. Tính sẵn sàng của nó là bao nhiêu, và hai thay đổi thiết
  kế nào nâng nó lên mà không đụng tới các phụ thuộc?
\end{enumerate}
\end{quiz}
